% ------------------------------------------------------------------------
% file `figures_planes-exercise-8-solution-body.tex'
%   in folder `xsim/'
%
%     solution of type `exercise' with id `8'
%
% generated by the `solution' environment of the
%   `xsim' package v0.21 (2022/02/12)
% from source `figures_planes' on 2026/08/09 on line 456
% ------------------------------------------------------------------------
\begin{enumerate}
    \item Pour placer le milieu $M$, aligne le grand bord gradué sur le segment $[AB]$ de manière à ce que le point $A$ soit sur la graduation $3\text{ cm}$ (à gauche) et le point $B$ soit sur la graduation $3\text{ cm}$ (à droite). Le point situé en face du \textbf{zéro central} est le milieu $M$ du segment $[AB]$. Ainsi, $AM = MB = 3\text{ cm}$.
    \item En maintenant le zéro central sur $M$, superpose la \textbf{ligne d'orthogonalité} sur le segment $[AB]$ et trace la droite $d$ le long du bord gradué. La droite $d$ passe par le milieu de $[AB]$ et lui est perpendiculaire : c'est la \textbf{médiatrice} du segment $[AB]$.
\end{enumerate}
